% Einheit 1 von 4 – Parameterform einer Ebene (bank/ebenen/e1.jsonl, zusammenbau v0.1)
%% TODO Zeitmarke Einheit 1: Marken-Zeile nicht lesbar
%% TODO Zweigzeile Teil 1: was hier gelernt wird (Satz in Schülersprache aus dem Einheitstitel) – Einheit 1
\einheitenkopf[e1]{Einheit 1 von 4 · Parameterform einer Ebene}
\zweigzeile{Abitur GK}
\setcounter{aufgabe}{9}

%% TODO Titel: Ich-kann-Satz fehlt in der Bank (Platzhalter: Kettenname) – Nr. 10
%% TODO Anweisung: ein Satz über der ersten Teilaufgabe fehlt in der Bank – Nr. 10
\begin{aufgabe}{Parameterform}
%% TODO \rechenplatz{n}: Schrittzahl des Grundfalls fehlt in der Bank (nur bei fester Schreibform, 2.3 b)
\begin{teile}
\teil Eine Ebene soll durch $A(1 | 3 | 0)$, $B(2 | 0 | 4)$ und $C(5 | 1 | 1)$ gehen. Was ist gegeben? \\ \kreuz{drei Punkte} \\ \kreuz{ein Punkt und zwei Richtungen} \\ \kreuz{eine Gerade und ein Punkt} \\ \kreuz{zwei Geraden}
\teil Punkt $(2 | 1 | 3)$, Spannvektoren $(1 | 0 | 4)$ und $(0 | 3 | 1)$: Parametergleichung der Ebene $E$?
\teil Die Punkte $A(1 | 0 | 3)$, $B(3 | 2 | 4)$, $C(5 | 1 | 2)$ und $D(3 | -1 | 1)$ sind in dieser Reihenfolge Eckpunkte eines Quadrats, das in der Ebene $E$ liegt. Gib eine Parametergleichung von $E$ an. \hfill (Abitur 2018 LK)
\teil Die Ebene $E$ enthält die Gerade $g\colon \vec x = (1 | 2 | 0) + t \cdot (1 | 1 | 2)$, $t \in \mathbb{R}$, und den Punkt $P(3 | 1 | 4)$. Parametergleichung von $E$?
\teil Ein Zeltdach hat vier rautenförmige Dachflächen; eine davon hat die Ecken $C(6 | 6 | 0)$, $F(6 | 3 | 3)$, $G(3 | 6 | 3)$ und die Spitze $S$ (1 LE = 1 m). Gib eine Parametergleichung der Ebene $L$ durch $C$, $F$ und $G$ an und weise nach, dass $S(3 | 3 | 6)$ in $L$ liegt. \hfill (Abitur 2022 GK)
\teil Beschreiben $E_1\colon \vec x = (1 | 1 | 0) + r \cdot (1 | 0 | 1) + s \cdot (0 | 1 | 1)$ und $E_2\colon \vec x = (1 | 1 | 0) + r \cdot (2 | 0 | 2) + s \cdot (0 | -3 | -3)$, $r, s \in \mathbb{R}$, dieselbe Ebene?
\teil Gegeben sind $A(2 | -1 | 1)$, $B(1 | 3 | 4)$ und der Punkt $C$ mit $\overrightarrow{OC} = 3 \cdot \overrightarrow{OA}$; $O$ ist der Koordinatenursprung. Begründe, dass es genau eine Ebene gibt, die $O$, $A$, $B$ und $C$ enthält. \hfill (Abitur 2018 GK)
\end{teile}
\end{aufgabe}

%% TODO Titel: Ich-kann-Satz fehlt in der Bank (Platzhalter: Kettenname) – Nr. 11
%% TODO Anweisung: ein Satz über der ersten Teilaufgabe fehlt in der Bank – Nr. 11
\begin{aufgabe}{Parameterform – Fehler finden, Begründen, Anwendung, Darstellungswechsel}
\begin{teile}
\teil Tim soll eine Parametergleichung der Ebene des Rechtecks $ABCD$ mit $A(1 | 2 | 1)$, $B(7 | 2 | 1)$, $C(7 | 5 | 5)$ und $D(1 | 5 | 5)$ angeben. Er schreibt $\vec x = (1 | 2 | 1) + r \cdot \overrightarrow{AB} + s \cdot \overrightarrow{DC}$ mit $\overrightarrow{AB} = (6 | 0 | 0)$ und $\overrightarrow{DC} = (6 | 0 | 0)$. Finde den Fehler.
\teil Begründe, warum eine Ebene unendlich viele verschiedene Parametergleichungen hat.
\teil Eine rechteckige Solarplatte hat die Ecken $A(0 | 0 | 3)$, $B(2 | 0 | 3)$, $C(2 | 1{,}5 | 5)$ und $D(0 | 1{,}5 | 5)$ (1 LE = 1 m). Ein Sensor soll bei $S(1 | 0{,}75 | 4)$ in der Plattenebene sitzen. Sitzt er dort?
\teil Gib drei Punkte der Ebene $E\colon \vec x = (1 | 2 | 0) + r \cdot (3 | 0 | 1) + s \cdot (0 | 1 | 2)$, $r, s \in \mathbb{R}$, an, die nicht auf einer Geraden liegen.
\end{teile}
\end{aufgabe}
